Krátké úvodní poznámky
\begin{itemize}
  \item Tato část nabízí praktický „rychlý start“ pro pedagoga, který chce poprvé vyzkoušet generativní AI v přípravě nebo v hodině. Obsahuje minimální bezpečnostní kroky, doporučený postup pilotu a konkrétní ukázku (Microsoft 365 Copilot + Forms). Informace o dostupnosti školních verzí Copilota jsou uvedeny níže \cite{kb_ai_101_tipu_pdf_04e4a115_page_3_1,kb_ai_101_tipu_pdf_04e4a115_page_36_1}.
  \item Poznámka: Seznam konkrétních služeb (např. ChatGPT, Claude, Gemini) uvádíme jako orientační — tyto stručné zmínky jsou označeny jako general background knowledge a nevyjadřují doporučení konkrétního produktu v dané instituci (general background knowledge).
\end{itemize}

Rychlý 6‑krokový workflow pro první použití
\begin{enumerate}
  \item Definujte cíl a rozsah pilotu: konkretizujte předmět, typ úlohy (např. generování cvičných testů, návrh slidů, krátká formativní zpětná vazba) a délku pilotu. Zároveň rozhodněte, které typy dat nebudou nikdy odesílány do externích služeb (osobní údaje, studentské práce s identifikátory apod.) \cite{kb_malinka_chatgpt_ve_vyuce_dva_roky_pote_2024_pdf_90bc3aaf_page_12_1}.
  \item Proveďte zjednodušené DPIA a datový plán: zdokumentujte, jaký datový tok bude v pilotu (co se posílá, kde se logy ukládají, kdo má k nim přístup) a uložte tuto dokumentaci jako součást pilotu \cite{kb_malinka_chatgpt_ve_vyuce_dva_roky_pote_2024_pdf_90bc3aaf_page_12_1}.
  \item Zvolte nástroj podle potřeby a licence: ověřte, zda vaše fakulta/univerzita nabízí enterprise nebo školní variantu (např. Office 365 pro školy/A1 obsahuje webovou verzi Copilot Chat; plná integrace Copilota do Word/PowerPoint/Forms je k dispozici v placené verzi Microsoft 365 Copilot) \cite{kb_ai_101_tipu_pdf_04e4a115_page_3_1}. Ujistěte se, že licence a DPA jsou v souladu s požadavky vašeho pracoviště.
  \item Začněte malým dvoufázovým pilotem (bez AI → s AI): porovnejte kontrolní skupinu bez AI a skupinu s AI, zaznamenejte artefakty a procesy (drafty, interakce s AI) pro pozdější vyhodnocení \cite{kb_malinka_chatgpt_ve_vyuce_dva_roky_pote_2024_pdf_90bc3aaf_page_12_1}.
  \item Nastavte pravidla „human‑in‑the‑loop“ a transparentnost: určete, kdo finálně schvaluje výstupy AI, jak probíhá eskalace chyb, a informujte studenty o použití AI (text do sylabu/oznámení) \cite{kb_malinka_chatgpt_ve_vyuce_dva_roky_pote_2024_pdf_90bc3aaf_page_12_1}.
  \item Vyhodnoťte pilot podle předem definovaných metrik: pedagogická kvalita výstupů, časová úspora, bezpečnostní incidenty, dopad na akademickou integritu; na základě výsledků upravte rozsah a procesy.
\end{enumerate}

Konkrétní příklad: rychlé vytvoření kvízu v Microsoft Forms pomocí Copilota
\begin{itemize}
  \item Kontext: Microsoft 365 Copilot je dostupný ve více variantách; školní Office 365 může obsahovat webovou verzi Copilot Chat a placená varianta Copilot nabízí integrace přímo do aplikací (Word, PowerPoint, Forms apod.) \cite{kb_ai_101_tipu_pdf_04e4a115_page_3_1}.
  \item Praktický postup:
    \begin{enumerate}
      \item Ověřte, že vaše školní/ fakultní účetní nastavení zahrnuje Copilot Chat nebo Microsoft 365 Copilot (kontaktujte IT/elearning pro potvrzení licence).
      \item Otevřete Microsoft Forms, aktivujte Copilot (pokud je dostupný) a použijte jednoduchý prompt pro vytvoření testu.
      \item Upravený výsledek zkontrolujte a doplňte manuálně (správné odpovědi, jmenné formátování, úprava úrovně obtížnosti).
    \end{enumerate}
  \item Ukázkový prompt (převzatý z praktických příkladů): 
\begin{verbatim}
Vygeneruj mi test na hodinu zeměpisu na téma Austrálie, test bude obsahovat otázky na základní geografické oblasti, bude pro žáky 8. ročníku základní školy v České republice, test bude mít 10 otázek, 4 otázky budou typu ABC, 4 otázky typu pravda/nepravda a 2 otázky otevřené.
\end{verbatim}
Tento způsob generování vrátí sadu otázek včetně označení správných odpovědí a často i krátkých vysvětlení — vždy ale prověřte přesnost a přiměřenost pro váš kontext před použitím v hodnocení \cite{kb_ai_101_tipu_pdf_04e4a115_page_36_1}.
\end{itemize}

Krátký bezpečnostní checklist před prvním použitím v hodině
\begin{itemize}
  \item Máte záznam DPIA nebo minimální datový plán? \cite{kb_malinka_chatgpt_ve_vyuce_dva_roky_pote_2024_pdf_90bc3aaf_page_12_1}
  \item Jsou citlivá data (jména, osobní identifikátory, studentské nahrávky) vyloučena nebo pseudonymizována? \cite{kb_malinka_chatgpt_ve_vyuce_dva_roky_pote_2024_pdf_90bc3aaf_page_12_1}
  \item Je jasně určeno, kdo finálně validuje výstupy AI (human‑in‑the‑loop)? \cite{kb_malinka_chatgpt_ve_vyuce_dva_roky_pote_2024_pdf_90bc3aaf_page_12_1}
  \item Máte potvrzení licence / dostupnosti Copilota nebo jiného nástroje u interní IT podpory (pokud plánujete používat školní Office 365/Copilot)? \cite{kb_ai_101_tipu_pdf_04e4a115_page_3_1}
  \item Plánujete krátké vyhodnocení pilotu (srovnání bez AI vs. s AI) a sběr zpětné vazby od studentů?
\end{itemize}

Drobné tipy pro rychlé nasazení
\begin{itemize}
  \item Začněte s nehodnoticím obsahem (příprava cvičných listů, návrhy slidů), kde případné chyby lze rychle opravit — to sníží riziko dopadu na studenty (general background knowledge).
  \item Při generování testů nebo úloh vždy proveďte rychlou kontrolu faktů a jazykové úpravy před sdílením se studenty; u automaticky generovaných řešení doplňte pedagogické komentáře.
  \item Pokud máte přístup k Microsoft 365 Copilot v rámci školního Office, využijte integrace s Forms pro rychlou tvorbu kvízů a s Word/PowerPoint pro převod sylabu do slidů — ale vždy zachovejte lidskou kontrolu výsledků \cite{kb_ai_101_tipu_pdf_04e4a115_page_3_1,kb_ai_101_tipu_pdf_04e4a115_page_36_1}.
\end{itemize}

Odkazy na další kroky v příručce
\begin{itemize}
  \item Pro podrobnější postupy, DPIA šablonu a texty do sylabu využijte přílohy 9.3–9.5 a online repozitář ai.vut.cz (living document) — doporučené šablony a checklisty jsou připravené pro rychlé zkopírování a přizpůsobení.
\end{itemize}